\noindent
The U.S. New Markets Tax Credit (NMTC), enacted in 2000, subsidizes private
investors with a 39\% federal tax credit conditional on routing capital
through certified Community Development Entities (CDEs) into projects in
low-income census tracts. Over FY2001--FY2022 the program deployed
\$66.6 billion of public credit into 8{,}024 projects worth \$120.9 billion,
implying an aggregate private-mobilization ratio of \$0.82 per federal
dollar. Non-metropolitan tracts received 19.6\% of dollars---effectively the
20\% statutory minimum---but project-level leverage in non-metro tracts is
systematically lower (median 1.07$\times$ versus 1.19$\times$ metro). This
paper decomposes the rural mobilization penalty using project-level
fixed-effects regressions on the CDFI Fund's public release. The naive
estimate is a precise $-0.262$ rural penalty in mean leverage. Year,
project-type, and state fixed effects shrink it to $-0.172$. Adding CDE
fixed effects moves the coefficient to $-0.047$; the mean regression is
imprecise (cluster-robust 95\% CI $[-0.245, +0.151]$), but the
median-regression analog is a precise $-0.001$ (SE $0.008$), and within
intermediary both the extensive margin (any private mobilization) and the
intensive margin are null. An order-invariant Gelbach decomposition
attributes $-0.185$ of the raw gap to CDE identity (CDE-cluster bootstrap
95\% CI $[-0.347, -0.023]$; 86\% of the explained movement), against
$-0.041$ for project type. The aggregate rural leverage gap is therefore
largely a between-CDE composition fact, not a within-CDE deployment
penalty. A cross-CDE bunching test on the 20\% non-metropolitan statutory
share finds no excess mass at the cumulative-deployment level (excess mass
$\hat B = -0.0006$; cluster-bootstrap 95\% CI $[-0.014, +0.013]$),
consistent with the mandate binding at the allocation-award level rather
than the deployment level. The findings reframe the rural mobilization
debate: the operative policy levers concern intermediary selection and
capacity, not market-structure redesign.
